Clasifique las ecuaciones siguientes como ordinarias o en derivadas parciales; indique su orden y su grado respecto de las derivadas en la forma escrita. Para las EDO, indique si son lineales. Los nombres entre paréntesis identifican contextos de origen y no requieren formular un modelo. Considere parámetros constantes; en el caso de difusión, $k>0$.

\begin{enumerate}[label=(\alph*),ref=\alph*,leftmargin=*]
\item $5x''(t)+2x'(t)+9x(t)=2\cos(3t)$ (vibraciones y circuitos).
\item $1+(y')^2=C$, con $C$ constante (cálculo de variaciones).
\item $8y^{(4)}=x(1-x)$ (deflexión de vigas).
\item $u_{xx}+u_{yy}=0$ (potencial y equilibrio térmico).
\item $y'=\dfrac{y(2-3x)}{x(1+3y)}$ (competencia de especies).
\item $-\pi(y\tan\alpha)^2y'=12\sqrt{2gy}$, con $g>0$ y $\tan\alpha$ finita y no nula (vaciado de un recipiente).
\item $u_t=k u_{xx}$, $k>0$ (difusión y conducción de calor).
\end{enumerate}
